% Standalone figure: run make in this directory to build PDF and SVG.
\documentclass[11pt,tikz,border=0pt]{standalone}
\usepackage{fontspec}
\IfFontExistsTF{Helvetica Neue Light}{
  \setsansfont{Helvetica Neue Light}[BoldFont={Helvetica Neue Medium}]
  \newfontfamily\vnmedium{Helvetica Neue Medium}
}{
  \setsansfont{TeX Gyre Heros}
  \newcommand{\vnmedium}{\bfseries}
}
\renewcommand{\familydefault}{\sfdefault}
\usepackage{amsmath,amssymb,mathtools}
\usepackage{unicode-math}
\setmathfont{latinmodern-math.otf}
\usetikzlibrary{arrows.meta,calc,positioning}
\definecolor{vnslideink}{HTML}{333333}
\definecolor{vnslidebody}{HTML}{5E5E5E}
\definecolor{vnaccent}{HTML}{FF644E}
\definecolor{vnslidelink}{HTML}{579ACA}
\pagecolor{white}

\begin{document}
% The panel export omits the title already supplied by the slide.
\begingroup
\definecolor{cgAccent}{HTML}{EF4136}
\definecolor{cgEdge}{HTML}{9A9A9A}
\definecolor{cgInk}{HTML}{333333}
\definecolor{cgBody}{HTML}{555555}
\definecolor{cgRule}{HTML}{D7D7D7}
\begin{tikzpicture}[x=1cm,y=1cm,
  font=\sffamily\fontsize{11}{13}\selectfont,text=cgInk,
  cgvertex/.style={circle,draw=cgAccent,fill=white,line width=1.1pt,
    minimum size=7mm,inner sep=0pt,font=\sffamily\bfseries\fontsize{11}{13}\selectfont},
  cglabel/.style={anchor=west,text=cgBody,font=\sffamily\fontsize{10}{12}\selectfont},
  cgequation/.style={anchor=west,font=\fontsize{13}{16}\selectfont}]
  \ifdefined\CompleteGraphOmitHeading
    \fill[white] (-0.5,-0.95) rectangle (11.1,4.25);
    \path[use as bounding box] (-0.5,-0.95) rectangle (11.1,4.25);
  \else
    \fill[white] (-0.5,-0.95) rectangle (11.1,5.55);
    \node[font=\sffamily\bfseries\fontsize{16}{19}\selectfont] at (5.3,5.2)
      {Complete graph $K_5$};
    \node[text=cgBody,font=\sffamily\fontsize{11}{13}\selectfont] at (5.3,4.7)
      {Every pair of distinct vertices is joined by an edge};
    \path[use as bounding box] (-0.5,-0.95) rectangle (11.1,5.55);
  \fi

  % Draw the six other edges and the four incident edges before the vertices.
  \coordinate (c1) at (2,3.8);
  \coordinate (c2) at (4,2.25);
  \coordinate (c3) at (3.25,0);
  \coordinate (c4) at (0.75,0);
  \coordinate (c5) at (0,2.25);
  \foreach \i/\j in {2/3,2/4,2/5,3/4,3/5,4/5}
    {\draw[draw=cgEdge,line width=0.85pt] (c\i)--(c\j);}
  \foreach \j in {2,3,4,5}
    {\draw[draw=cgAccent,line width=1.45pt] (c1)--(c\j);}
  \foreach \i in {2,3,4,5}
    {\node[cgvertex] at (c\i) {\i};}
  \node[cgvertex,fill=cgAccent!7,line width=1.6pt] at (c1) {1};

  \node[anchor=west,font=\sffamily\bfseries\fontsize{11}{13}\selectfont]
    at (5.4,3.9) {What the picture shows};
  \draw[draw=cgRule,line width=0.5pt] (5.4,3.55)--(10.9,3.55);
  \node[cglabel] at (5.4,3.25) {Degree of every vertex};
  \node[cgequation] at (5.4,2.8) {$\displaystyle\deg(v)=5-1=4$};
  \node[cglabel] at (5.4,2.15) {Number of edges};
  \node[cgequation] at (5.4,1.2) {$\displaystyle|\mathcal E|=\frac{5(4)}{2}=10$};
  \node[cglabel] at (5.4,0.35) {Fraction of possible edges};
  \node[cgequation] at (5.4,-0.1) {$\displaystyle\rho(K_5)=1$};
  \draw[draw=cgAccent,line width=1.45pt] (5.4,-0.7)--(6.0,-0.7);
  \node[anchor=west,text=cgBody,font=\sffamily\fontsize{9}{11}\selectfont]
    at (6.15,-0.7) {edges incident to vertex 1};
\end{tikzpicture}
\endgroup

\end{document}
